% ----------------------------------------------------------------------------
%  12. Comparative Analysis
% ----------------------------------------------------------------------------
\section{Comparative Analysis}
\label{sec:comparative}

\subsection{Greedy vs.\ Hungarian, head-to-head}

\begin{table}[H]
\centering
\small
\begin{tabular}{l r r r}
\toprule
\textbf{Axis} & \textbf{Greedy} & \textbf{Hungarian} & \textbf{Difference} \\
\midrule
Makespan (min)             & 9{,}260 & 9{,}470 & $+210$ \\
Total overtime (min)       & 20{,}210 & 20{,}210 & 0 \\
Mean critical finish (min) & 1{,}795 & 3{,}318 & $\bm{+85\%}$ \\
Worst critical finish (min)& 2{,}415 & 9{,}470 & $\bm{4\times}$ \\
Wall-clock                 & 16~ms   & 84~ms   & $+428\%$ \\
\bottomrule
\end{tabular}
\end{table}

\subsection{Why the textbook is wrong here}

\begin{itemize}
  \item Hungarian solves the ``best one-to-one match for this
        single tick'' problem perfectly.
  \item But many critical tasks arrive together with the same
        priority. Hungarian places one per worker per tick;
        the rest are pushed to the next tick.
  \item Greedy stacks all of them on free workers in the same
        tick. The deferral never happens.
  \item The textbook claim --- ``Hungarian is optimal'' --- is true
        only at the single-tick level. Across many ticks, Greedy's
        stacking is what the patients actually experience.
\end{itemize}

\subsection{Fairness: who carries the load?}

\begin{itemize}
  \item Within each role, Greedy already balances the work
        (within-role spread $\approx 7$~min).
  \item ``Fair-mode'' rebalancing only helps when workers
        \emph{start} unequal. They don't on this ward.
  \item The big imbalance is \emph{across} roles, not within them
        --- and that's a staffing problem, not a matching problem.
\end{itemize}

\subsection{Why no algorithm can fix this}

\begin{itemize}
  \item Total work needed: $\sim$25{,}000~min.
  \item Total available shift time: 5{,}040~min.
  \item Workload is $\sim$5$\times$ over capacity.
  \item Two nurses end the day at 9{,}260 / 9{,}255 min worked.
        Three interns end at 240~min each. A 39$\times$ ratio.
  \item No ordering or matching change can move that.
\end{itemize}

\begin{insightbox}
\textbf{The one number that matters.}\;
20{,}210~minutes of overtime show up under \emph{every} strategy.
The fix is more nurses, not a smarter scheduler.
\end{insightbox}
